\scanpage{121}
\mspage{115}
% archive sheet number: 120

les $a\in A$ sont ``\uline{commensurables dans leur} ensemble'' (\uline{NB} ce qui est plus fort que d'exiger qu'ils le soient $2:2$), si $A\subset\Sigma'$, $\Sigma'$ \unsure{un|une} \deleted{\unsure{p.}} ``sous-spatios'' (stable par $\operatorname{Sup}$, $\unsure{\ell|\complement}$ donc par $\operatorname{Inf}$) qui \uline{est une alg. de Boole}.

Ceci posé, étant donné un diagramme
\[
A_0\ll A_1\ll\cdots\ll A_n,
\]
nous dirons qu'il est commensurable, si $A=\bigcup_{1\leq i\leq n}A_i$ l'est.

\deleted{On pourrait encore \unsure{l'exp}}

\uline{Lemme 3} Pour tout $1\leq i\leq n$, et $a_i\in A_i$, considérons $a_i''$ défini comme dans lemme 2, pour $A_{i-1}\ll A_i$. Soit
\[
A_i''=\{a_i''\mid a_i\in A_i\text{ t.q. }\unsure{c_i''|a_i''}\ne\deleted{\unsure{\varnothing_\Sigma|\varnothing}}0_\Sigma\}.
\]
\deleted{Soit} On pose aussi
\[
A_0''=A_0,
\]
\deleted{$\unsure{A}$}
\[
A=\bigcup_{0\leq i\leq n}A_{\unsure{j|i}}''.
\]
Alors les $A_i''$ sont formés d'éléments \inserted{non nuls} mut. disj, et ils sont mut. disjoints entre eux,
